\section{The route-selection question}
Turn 2 constructed the complete canonical shifted-factor graph and found closed
components on which both original selectors fail. Those components contained
large, changing obstruction quotients. The narrower possibility remained that
\emph{the smallest active quotient, of order six, could not persist around a
closed component}. This tranche disproves that narrower assertion as well.

\begin{theorem}[Closed primitive sextic triangle]\label{thm:triangle}
Let
\[
 p=113946012068401\equiv121\pmod{840}.
\]
This is a prime. Its canonical nonresidue-factor graph has the closed component
\[
 \boxed{31\longrightarrow223\longrightarrow307\longrightarrow31.} \tag{1}
\]
At each vertex the original exterior and middle selectors are empty, the actual
full unit group and coefficient envelope agree, their actual support stabilizer
has index six, and the inactive original divisor block fills that stabilizer.
The exceptional successor prime is simple and is the only outgoing nonresidue
prime. Every vertex has the first phase $pK=rK$; every neighbor-square relation
has an actual sixth-root lift. The three complete boxes have respectively
$45,405,405$ original vectors.
\end{theorem}
The statement is a counterexample to a proposed \emph{proof implication}, not to
Erd\H{o}s--Straus (ES). Section~\ref{sec:positive} returns a positive integral
solution at this same prime. No least-prime assertion or complete scan of the
large prime's entire graph is made.

The universal source graph is due to Bryan \cite[\S\S1--7]{BryanGraph}.
The primitive sextic normal form and the neighbor-square recurrence are also
antecedents \cite[\S\S3--8]{BryanSix}, \cite[\S\S2--4]{BryanChain}.
The latter gives two consecutive primitive failures at $p=808369$, then a larger
index. Turn 2 supplied closed components containing a primitive failure
\cite[\S\S7--8]{Turn2}. The new certificate keeps \emph{every} vertex primitive
and closes in three edges. A content-level search found no matching certificate;
this is not a historical-priority determination.

\section{The original source and a certificate criterion}
For a hard prime $p$, the vertices are primes $q<p$ with $(q/p)=-1$.
Here all three vertices satisfy $q\equiv3\pmod4$, so
\[
 a_q=(p+q)/4,\qquad R_q=q,\qquad
 a_q=\prod_t t^{e_t},\qquad
 \mathcal B_q=\prod_t[-e_t,e_t]_{\mathbb Z}.                 \tag{2}
\]
For each original vector retain
\[
 v(\beta)=\prod_t t^{\beta_t}\pmod q,\quad
 u(\beta)=\prod_t t^{e_t+\beta_t}\mid a_q^2,\quad
 \mu_q(x)=\#\{\beta:v(\beta)=x\},\quad W_q=\operatorname{supp}\mu_q. \tag{3}
\]
The inverse is $\beta_t=v_t(u)-e_t$. In particular $\mu_q(x)=\mu_q(x^{-1})$.
The canonical gates and their centered targets are
\[
 M:q\mid u+a_q\ \longleftrightarrow\ v=-1,\qquad
 E:q\mid4u+1\ \longleftrightarrow\ v=-p^{-1}.              \tag{4}
\]
The reflected exterior target $-p$ requires the retained map $\beta\mapsto-\beta$;
it is not a third canonical channel. All prime-power multiplicities in (2) remain.

The full edge set consists of \emph{all} primes $r\mid a_q$ with $(r/p)=-1$, with
weight $v_r(a_q)$. An edge is not chosen by a heuristic rule. In this tranche the
complete factorization proves that every vertex has exactly one edge.

Here is a direct sufficient test for the claimed failure structure. Set
$G_q=(\mathbb Z/q\mathbb Z)^\times$, $K_q=G_q^6$, and suppose the actual integer
factorization is $a_q=B_qr$, with $r\nmid B_q$, all factors of $B_q$ in $K_q$,
\[
 W_q(B_q)=K_q,\qquad \operatorname{ord}(rK_q)=6,\qquad4\in K_q. \tag{5}
\]
The phase is then fixed by the original equation:
\[
 p\equiv4B_qr\pmod q\quad\Longrightarrow\quad pK_q=rK_q.
\]
Orient $rK_q$ as class $1$ in $C_6$. The complete support and targets are
\[
 W_q/K_q=\{0,1,5\},\qquad
 (M,E_+,E_-)=(3,2,4).                                    \tag{6}
\]
Thus both channels fail. The set $\{0,1,5\}$ has trivial stabilizer in $C_6$;
therefore the stabilizer of $W_q$ in the full unit group is exactly $K_q$, not a
guessed power-residue group. Since $W_q(B_q)=K_q$ and $r$ is an original factor,
$\Gamma_q=\langle-1,2,t:t\mid a_q\rangle$ contains $K_q$ and its quotient generator,
so $\Gamma_q=G_q$. Both indices are six.

Every prime factor $t\mid a_q$ satisfies $(t/p)=(t/q)$. For odd $t$, use
$p\equiv-q\pmod t$, $p\equiv1\pmod4$, $q\equiv3\pmod4$ and quadratic reciprocity;
for $t=2$ use the supplementary law. The reciprocity input is classical
\cite[\S6.4]{Hatcher}. The inactive factors are squares modulo $q$, whereas $r$
is not, since $q-1$ has exactly one factor of two. This proves that $r$ is the
only outgoing nonresidue factor. The code also checks the characters directly.

\paragraph{Minimality of the structural sizes.}
At any external nonresidue prime residual $q\equiv3\pmod4$, a failed stabilizer
cannot contain $-1$, so it has odd order. Its index is twice an odd integer.
Index two would make the stabilizer the squares, containing the exterior target
$-p^{-1}$, since $(p/q)=-1$. Thus six is the first possible index.
There are no self-edges because $q\nmid a_q$. A two-cycle of $3\pmod4$ primes
would require both reciprocal Legendre symbols to be $-1$, contrary to quadratic
reciprocity. Hence three is the shortest possible cycle in this prime-residual
sector. Theorem~\ref{thm:triangle} realizes both minimal structural sizes, not a
minimum value of $p$.

\section{Complete integer and coefficient certificate}\label{sec:certificate}
The three source factorizations are
\begin{align*}
 a_{31}&=28486503017108=223\,B_{31},& B_{31}&=2^2\cdot31935541499,\\
 a_{223}&=28486503017156=307\,B_{223},& B_{223}&=2^2\cdot17\cdot41\cdot33281891,\\
 a_{307}&=28486503017177=31\,B_{307},& B_{307}&=19^2\cdot269\cdot937\cdot10099.
\end{align*}
All displayed factors are prime. Complete trial division certifies the factors.
For the large parameter,
\[
 p-1=2^4\cdot3^3\cdot5^2\cdot10550556673.                 \tag{7}
\]
The last factor is prime, and base $11$ satisfies $11^{p-1}\equiv1\pmod p$.
The residues of $11^{(p-1)/\ell}$ for the four distinct primes $\ell$ in (7) are
\[
 113946012068400,\quad68851444066,\quad71881191975328,\quad61379549739175.
\tag{8}
\]
Each residue minus one is coprime to $p$. If a prime $d$ divided $p$, the order
of $11\bmod d$ would therefore contain every full prime power in $p-1$.
Thus $p-1\mid d-1$, forcing $d\ge p$ and proving $p$ prime.
The separate implementation additionally tries every odd divisor through
$\lfloor\sqrt p\rfloor=10674549$.

The remaining certificate is small. Use kernel generators $b_q$ and centered
logarithmic factors as follows; each ordered pair is (logarithm, original exponent).
\[
\begin{array}{c|r|r|l|r}
q&|K_q|&b_q&\text{inactive centered factors in }C_{|K_q|}&\tau(B_q^2)\\\hline
31&5&2&(1,2),(0,1)&15\\
223&37&2&(1,2),(23,1),(11,1),(8,1)&135\\
307&51&4&(19,2),(45,1),(2,1),(28,1)&135
\end{array}                                                  \tag{9}
\]
For example, $33281891\equiv2^8\pmod{223}$ and $10099\equiv4^{28}\pmod{307}$.
The displayed generators have exact orders $5,37,51$. Direct exponentiation
checks every factor and order. The complete integer coefficient products are
\[
 \prod_{(d,e)\text{ in row }q}\sum_{j=-e}^{e}X^{dj}
 \equiv\sum_{j=0}^{|K_q|-1}c_{q,j}X^j\pmod{X^{|K_q|}-1},       \tag{10}
\]
where the following are the \emph{entire} coefficient vectors:
\begingroup\small
\begin{align*}
c_{31}={}&(3,3,3,3,3),\\
c_{223}={}&(1,3,4,4,5,5,4,3,2,3,4,3,4,5,3,3,4,4,4,\\
           &\quad4,4,4,3,3,5,4,3,4,3,2,3,4,5,5,4,4,3),\\
c_{307}={}&(3,1,4,1,4,2,3,3,3,3,2,4,2,4,1,5,1,5,1,4,1,4,1,3,2,2,\\
           &\quad2,2,3,1,4,1,4,1,5,1,5,1,4,2,4,2,3,3,3,3,2,4,1,4,1).
\end{align*}
\endgroup
These polynomial identities can be checked by finite multiplication. Every
coefficient is positive, proving $W(B_q)=K_q$ with its original multiplicities.
The successor orders in the full groups are $6,222,306$, respectively; their
classes all have order six. Each $4\in K_q$. Thus (5)--(6) apply.

The exact fine targets are
\[
\begin{array}{c|rrr|r|r}
q&M&E_+&E_-&|W_q|&|\mathcal B_q|\\\hline
31&30&10&28&15&45\\
223&222&36&31&111&405\\
307&306&201&223&153&405.
\end{array}                                                   \tag{11}
\]
All targets are in the three unoccupied quotient classes. The complete $855$
original vectors, their divisors $u$, normalized $h,r,s$, channel tags, raw
quotients and ordered rational denominator triples are serialized in JSON.
Every candidate rational identity equals $4/p$; every one fails the required
integral gate in both channels. No original exponent vector is omitted.
The unique outgoing edges in (1) follow from the complete factorizations and
character checks. This proves Theorem~\ref{thm:triangle}.

\section{The odd component and the exact lost correlation}\label{sec:mixed}
The following calculation applies to \emph{any} primitive sextic failure with
one simple active prime and saturated inactive factor $B$. Write $T=\tau(B^2)$;
$T$ is odd. In the quotient oriented by the active prime, its original counting
polynomial is
\[
 \overline\mu=T(1+g+g^{-1})\in\mathbb Z[C_6].              \tag{12}
\]
Here the bar denotes the integer quotient observation, not reduction modulo two.
The exact CRT coordinate map is $j\mapsto(j\bmod2,j\bmod3)$, inverse
$(a,b)\mapsto3a+4b\pmod6$. It turns (12) into the joint matrix
\[
 \boxed{C=T\begin{pmatrix}1&0&0\\0&1&1\end{pmatrix}.}     \tag{13}
\]
Its row sums are $(T,2T)$ and its column sums are $(T,T,T)$. In particular, the
\emph{odd cubic marginal is exactly uniform and positive}, while all three joint
target cells are absent. The data do not lack quadratic or cubic residues
separately; they lack the required joint pairings.

\begin{proposition}[Full integer marginal fibres]
Every nonnegative integer matrix with those row and column sums is uniquely
\[
 \begin{pmatrix}T-u-v&u&v\\u+v&T-u&T-v\end{pmatrix},
 \qquad u,v\ge0,\quad u+v\le T.                          \tag{14}
\]
There are $(T+1)(T+2)/2$ such matrices. Imposing the original inversion symmetry
reduces the entire fibre to
\[
 \boxed{C(w)=\begin{pmatrix}T-2w&w&w\\2w&T-w&T-w\end{pmatrix},
 \quad w\in\mathbb Z_{\ge0},\quad2w\le T.}                \tag{15}
\]
The sum of its two canonical target-cell coefficients is $3w$.
\end{proposition}
\begin{proof}
Read $u=C_{0,1}$ and $v=C_{0,2}$. The prescribed sums force every other entry in
(14), giving both the inverse and the nonnegativity inequalities. Inversion
exchanges cubic coordinates 1 and 2 without changing parity, so $u=v=w$.
The canonical E target is cell $(0,2)$ and the M target is $(1,0)$, giving $w+2w$. These are quotient-coset totals. They are not asserted to be
the exact fine exterior-plus-middle state count.
\end{proof}
These are fibres of the stated coefficient observation, not a claim that every
matrix in the fibre is the divisor distribution of an arithmetic shell. The
actual triangle realizes $w=0$ at all three vertices.

The integer kernel of taking row and column sums is generated by
\[
 \begin{pmatrix}1&-1&0\\-1&1&0\end{pmatrix},\qquad
 \begin{pmatrix}1&0&-1\\-1&0&1\end{pmatrix}.               \tag{16}
\]
Formula (14) supplies the complete inverse after its two kernel coordinates are
retained. Under inversion symmetry only the sum of those directions remains.
No independent-product replacement of the marginals is made.

There is also an exact form inside the characteristic-two coefficient algebra.
Write $g=yx$, with $y^2=1,x^3=1$, and put $\varepsilon=y+1$. Then
\[
 \mathbb F_2[C_6]\cong
 \mathbb F_2[\varepsilon]/\varepsilon^2\ \times\
 \mathbb F_4[\varepsilon]/\varepsilon^2.                  \tag{17}
\]
If $J=1+x+x^2$ and $E=x+x^2$, these are orthogonal idempotents with $J+E=1$.
Evaluation at $x=1$ and at a marked $\zeta\in\mathbb F_4$ satisfying
$\zeta^2+\zeta+1=0$ gives (17). Its inverse sends
$(a+b\varepsilon,c+d\varepsilon)$, with
$c=c_0+c_1\zeta,d=d_0+d_1\zeta$, to
\[
 (a+b\varepsilon)J+(c_0+c_1x)E+\varepsilon(d_0+d_1x)E.
\]
The actual primitive packet reduces to
\[
 \boxed{1+y(x+x^2)=J+\varepsilon E\longleftrightarrow(1,\varepsilon).} \tag{18}
\]
It is nonzero on the primitive cubic component but has square zero there.
The complete multiplication is verified on all $64^2$ pairs of coefficient
vectors, independently of the arithmetic counterexample.

Characteristic two is only an observation: $C(0)$ and $C(2)$ have the same
reduction when $T\ge4$, although their target-cell sums are 0 and 6. The full
integer $w$ and original fibres must remain.

Squaring the quotient packet is also not a permitted reuse of the original
factor budget:
\[
 (1+g+g^{-1})^2=3+2g+2g^{-1}+g^2+g^{-2}.                 \tag{19}
\]
The apparent E-target coefficients at $g^{\pm2}$ come exclusively from the two
active exponent pairs $(1,1)$ and $(-1,-1)$. Their sums $\pm2$ lie outside the
actual interval $[-1,1]$. The map of two source boxes by addition has its full
bounded-composition fibres, but its target is a \emph{doubled} exponent box.
It is not the original selector source. Across different vertices, even the
moduli and prime factors differ; quotient labels do not create an integral
intertwiner between the original boxes.

\section{Actual carries and cyclic arithmetic}
At vertex $q$ put $r=q_+$ and choose section $\sigma_q(j)=r^j$ for $0\le j<6$.
Every fine unit has unique coordinates $\sigma_q(j)k$, $k\in K_q$.
With $\delta_q=r^6$, the complete multiplication is
\[
 (i,k)(j,l)=((i+j)\bmod6,\ kl\delta_q^{\lfloor(i+j)/6\rfloor}). \tag{20}
\]
The cyclic carries are
\[
\begin{array}{c|r|r|r}
q&r\bmod q&\delta_q&\operatorname{ord}(\delta_q)\\\hline
31&6&1&1\\
223&84&8&37\\
307&31&144&51.
\end{array}                                                  \tag{21}
\]
The last extension $C_{306}\to C_6$ with kernel $C_{51}$ has no homomorphic
section: any lift of a quotient generator has order divisible by nine, so cannot
have sixth power one. Direct enumeration verifies the complete lift locus.
The other two group extensions split, but the marked section at 223 still has
its nontrivial carry. No direct-product substitution deletes (20).

If $\mu_B$ is the full inactive count, the fine quotient-array formula is
\[
 \nu(0,k)=\nu(1,k)=\mu_B(k),\quad
 \nu(5,k)=\mu_B(\delta_q k),\quad
 \nu(2,k)=\nu(3,k)=\nu(4,k)=0.                            \tag{22}
\]
Thus support saturation never requires uniform fine coefficients.

All three neighbor relations use the same sign:
\[
 q_-q_+^2\equiv u_q^6\pmod q.                            \tag{23}
\]
One exact choice is $(u_{31},u_{223},u_{307})=(3,38,20)$; the corresponding
kernel values are $(16,196,17)$. The complete six-element root fibres are in
JSON. Both neighbors have the same nontrivial cubic character at each vertex:
\[
 (q_-^{(q-1)/3},q_+^{(q-1)/3})=(25,25),(183,183),(17,17)
\]
in the respective fields. These are three separately marked field values, not
scalars multiplied across different finite fields.

The integer cycle itself has a full inverse. Writing $c_i=B_{q_i}$, let
\[
 D=64c_0c_1c_2,
 \quad T_i=1+4c_i+16c_ic_{i+1},\quad h=\gcd(T_0,T_1,T_2).
\]
The original equations give $(D-1)q_i=pT_i$ and hence
\[
 p=(D-1)/h,\qquad q_i=T_i/h,
 \quad h=6117772570470359649583.                          \tag{24}
\]
Every coefficient and cyclic starting label is retained.

The discovery progression, not an asserted all-failure progression, was
\[
 p(t)=1396761361+1782724440t,\qquad t=63916,
\]
with the three inactive affine values
\[
 1565876+1998570t,\quad1137428+1451730t,\quad11264207+14376810t.
\]
The edge equations hold identically there. The prime, inactive factorizations
and saturation at the specified $t$ are individually proved. No infinitude of
closed sextic components in this progression is inferred.

\section{A real ES return and the route-selection consequence}\label{sec:positive}
The counterexample prime has a separate original exterior state
\[
 R=3,\quad a=28486503017101=19\cdot127\cdot1949\cdot6057173,
 \quad u=1949,
\]
\[
 (h,r,s,\kappa)=(1949,1,14615958449,37986876008950).
\]
It returns the positive ordered denominators
\[
 \begin{split}
 x&=28486503017101,\\
 y&=1082113258039195768479053950,\\
 z&=8436154959673348101798680263550.
 \end{split}                                               \tag{25}
\]
Indeed $a=hrs$, $3\kappa=pr+s$, and the inherited exterior identity
\cite[\S2, Proposition 2.2]{ET} gives $4/p=1/x+1/y+1/z$ by direct multiplication.
Its ordered inverse is $u=a^2/(3y-pa)$. Primality of all factors and the rational
identity are checked in both implementations.

An actual additional port also repairs a vertex of the example. At $q=31$ the
square multiplier 9 gives
\[
 R'=279,\ a'=28486503017170,\ u'=1991,\quad
 (h,r,s,\lambda)=(1991,1,14307635870,51281849).             \tag{26}
\]
These give a genuine M return by $279\lambda=r+s$ and $a'=hrs$.
The multiplier is a square modulo $p$, satisfies the original size bound, and is
retained as a new source marking. This proves only this repair, not a universal
nine-port theorem. The unchanged original cycle remains closed in the canonical
graph.

There can be no strictly decreasing well-founded function of the original
marked vertices along \emph{every} forced edge: (1) returns to the identical
state after three steps. This excludes local descent on that graph; it does not
exclude a larger state space or a selective different arithmetic operation.
Nor does it refute the weaker claim that some other canonical component has a
hit. We have not enumerated the entire graph at this large prime.

The mandatory Turn 3 gate therefore closes the proposed canonical sextic-cycle
escape route. All the smallest odd-component data, nonuniform fine counts,
section corrections, target orientations, and integer cycle identities coexist
with failure. Turn 4 must derive a relation between \emph{different original
sources} or attack the complete original-shell signed count. Extending this
failed local-closure assertion to higher quotient orders would be invalid.
Equations (14)--(18) identify the missing mixed coefficient, not a positive
lower bound for it at every prime.

\section{Executable scope}
The standard-library main verifier checks all three factorizations, $855$
original vectors and both gate tests on each, both rational denominator returns
and their inverses, the full stabilizers, both full/effective indices, every
inactive coefficient, all carries, and all six-root fibres. Its primality proof
uses (7)--(8). It checks all $64^2$ products in (17), exact nonnegative marginal
fibres for $0\le T\le24$, and eight explicitly false transformations.
The separate implementation imports neither the main verifier nor predecessor
software. It uses direct divisors, Jacobi characters, literal support translates,
a separately expanded three-cycle inverse, and complete trial primality of $p$.
Run receipts and hashes distinguish finite verification from the written proofs.
No entire previous-tranche re-audit, remote publication without a receipt, new
prime-coverage range, independently reviewed proof, or Lean compilation is claimed.
